%!TeX root=../gardentop.tex
\chapter{»There Was Some One Crying—There Was!«}
\begin{figure}[t!]
\centering
\includegraphics[width=\linewidth]{may}
\end{figure}

	\lettrine[lines=6]{T}{he} next day the rain poured down in torrents again, and when Mary looked out of her window the moor was almost hidden by gray mist and cloud. There could be no going out to-day.

\zz
»What do you do in your cottage when it rains like this?« she asked Martha.

»Try to keep from under each other's feet mostly,« Martha answered. »Eh! there does seem a lot of us then. Mother's a good-tempered woman but she gets fair moithered. The biggest ones goes out in th' cow-shed and plays there. Dickon he doesn't mind th' wet. He goes out just th' same as if th' sun was shinin'. He says he sees things on rainy days as doesn't show when it's fair weather. He once found a little fox cub half drowned in its hole and he brought it home in th' bosom of his shirt to keep it warm. Its mother had been killed nearby an' th' hole was swum out an' th' rest o' th' litter was dead. He's got it at home now. He found a half-drowned young crow another time an' he brought it home, too, an' tamed it. It's named Soot because it's so black, an' it hops an' flies about with him everywhere.«

The time had come when Mary had forgotten to resent Martha's familiar talk. She had even begun to find it interesting and to be sorry when she stopped or went away. The stories she had been told by her Ayah when she lived in India had been quite unlike those Martha had to tell about the moorland cottage which held fourteen people who lived in four little rooms and never had quite enough to eat. The children seemed to tumble about and amuse themselves like a litter of rough, good-natured collie puppies. Mary was most attracted by the mother and Dickon. When Martha told stories of what »mother« said or did they always sounded comfortable.

»If I had a raven or a fox cub I could play with it,« said Mary. »But I have nothing.«

Martha looked perplexed.

»Can tha' knit?« she asked.

»No,« answered Mary.

»Can tha' sew?«

»No.«

»Can tha' read?«

»Yes.«

»Then why doesn't tha' read somethin', or learn a bit o' spellin'? Tha'st old enough to be learnin' thy book a good bit now.«

»I haven't any books,« said Mary. »Those I had were left in India.«

»That's a pity,« said Martha. »If Mrs~Medlock'd let thee go into th' library, there's thousands o' books there.«

Mary did not ask where the library was, because she was suddenly inspired by a new idea. She made up her mind to go and find it herself. She was not troubled about Mrs~Medlock. Mrs~Medlock seemed always to be in her comfortable housekeeper's sitting-room down-stairs. In this queer place one scarcely ever saw any one at all. In fact, there was no one to see but the servants, and when their master was away they lived a luxurious life below stairs, where there was a huge kitchen hung about with shining brass and pewter, and a large servants' hall where there were four or five abundant meals eaten every day, and where a great deal of lively romping went on when Mrs~Medlock was out of the way.

Mary's meals were served regularly, and Martha waited on her, but no one troubled themselves about her in the least. Mrs~Medlock came and looked at her every day or two, but no one inquired what she did or told her what to do. She supposed that perhaps this was the English way of treating children. In India she had always been attended by her Ayah, who had followed her about and waited on her, hand and foot. She had often been tired of her company. Now she was followed by nobody and was learning to dress herself because Martha looked as though she thought she was silly and stupid when she wanted to have things handed to her and put on.

»Hasn't tha' got good sense?« she said once, when Mary had stood waiting for her to put on her gloves for her. »Our Susan Ann is twice as sharp as thee an' she's only four year' old. Sometimes tha' looks fair soft in th' head.«

Mary had worn her contrary scowl for an hour after that, but it made her think several entirely new things.

She stood at the window for about ten minutes this morning after Martha had swept up the hearth for the last time and gone down-stairs. She was thinking over the new idea which had come to her when she heard of the library. She did not care very much about the library itself, because she had read very few books; but to hear of it brought back to her mind the hundred rooms with closed doors. She wondered if they were all really locked and what she would find if she could get into any of them. Were there a hundred really? Why shouldn't she go and see how many doors she could count? It would be something to do on this morning when she could not go out. She had never been taught to ask permission to do things, and she knew nothing at all about authority, so she would not have thought it necessary to ask Mrs~Medlock if she might walk about the house, even if she had seen her.

She opened the door of the room and went into the corridor, and then she began her wanderings. It was a long corridor and it branched into other corridors and it led her up short flights of steps which mounted to others again. There were doors and doors, and there were pictures on the walls. Sometimes they were pictures of dark, curious landscapes, but oftenest they were portraits of men and women in queer, grand costumes made of satin and velvet. She found herself in one long gallery whose walls were covered with these portraits. She had never thought there could be so many in any house. She walked slowly down this place and stared at the faces which also seemed to stare at her. She felt as if they were wondering what a little girl from India was doing in their house. Some were pictures of children—little girls in thick satin frocks which reached to their feet and stood out about them, and boys with puffed sleeves and lace collars and long hair, or with big ruffs around their necks. She always stopped to look at the children, and wonder what their names were, and where they had gone, and why they wore such odd clothes. There was a stiff, plain little girl rather like herself. She wore a green brocade dress and held a green parrot on her finger. Her eyes had a sharp, curious look.

»Where do you live now?« said Mary aloud to her. »I wish you were here.«

Surely no other little girl ever spent such a queer morning. It seemed as if there was no one in all the huge rambling house but her own small self, wandering about up-stairs and down, through narrow passages and wide ones, where it seemed to her that no one but herself had ever walked. Since so many rooms had been built, people must have lived in them, but it all seemed so empty that she could not quite believe it true.

It was not until she climbed to the second floor that she thought of turning the handle of a door. All the doors were shut, as Mrs~Medlock had said they were, but at last she put her hand on the handle of one of them and turned it. She was almost frightened for a moment when she felt that it turned without difficulty and that when she pushed upon the door itself it slowly and heavily opened. It was a massive door and opened into a big bedroom. There were embroidered hangings on the wall, and inlaid furniture such as she had seen in India stood about the room. A broad window with leaded panes looked out upon the moor; and over the mantel was another portrait of the stiff, plain little girl who seemed to stare at her more curiously than ever.

»Perhaps she slept here once,« said Mary. »She stares at me so that she makes me feel queer.«

After that she opened more doors and more. She saw so many rooms that she became quite tired and began to think that there must be a hundred, though she had not counted them. In all of them there were old pictures or old tapestries with strange scenes worked on them. There were curious pieces of furniture and curious ornaments in nearly all of them.

In one room, which looked like a lady's sitting-room, the hangings were all embroidered velvet, and in a cabinet were about a hundred little elephants made of ivory. They were of different sizes, and some had their mahouts or palanquins on their backs. Some were much bigger than the others and some were so tiny that they seemed only babies. Mary had seen carved ivory in India and she knew all about elephants. She opened the door of the cabinet and stood on a footstool and played with these for quite a long time. When she got tired she set the elephants in order and shut the door of the cabinet.

In all her wanderings through the long corridors and the empty rooms, she had seen nothing alive; but in this room she saw something. Just after she had closed the cabinet door she heard a tiny rustling sound. It made her jump and look around at the sofa by the fireplace, from which it seemed to come. In the corner of the sofa there was a cushion, and in the velvet which covered it there was a hole, and out of the hole peeped a tiny head with a pair of frightened eyes in it.

Mary crept softly across the room to look. The bright eyes belonged to a little gray mouse, and the mouse had eaten a hole into the cushion and made a comfortable nest there. Six baby mice were cuddled up asleep near her. If there was no one else alive in the hundred rooms there were seven mice who did not look lonely at all.

»If they wouldn't be so frightened I would take them back with me,« said Mary.

She had wandered about long enough to feel too tired to wander any farther, and she turned back. Two or three times she lost her way by turning down the wrong corridor and was obliged to ramble up and down until she found the right one; but at last she reached her own floor again, though she was some distance from her own room and did not know exactly where she was.

»I believe I have taken a wrong turning again,« she said, standing still at what seemed the end of a short passage with tapestry on the wall. »I don't know which way to go. How still everything is!«

It was while she was standing here and just after she had said this that the stillness was broken by a sound. It was another cry, but not quite like the one she had heard last night; it was only a short one, a fretful, childish whine muffled by passing through walls.

»It's nearer than it was,« said Mary, her heart beating rather faster. »And it is crying.«

She put her hand accidentally upon the tapestry near her, and then sprang back, feeling quite startled. The tapestry was the covering of a door which fell open and showed her that there was another part of the corridor behind it, and Mrs~Medlock was coming up it with her bunch of keys in her hand and a very cross look on her face.

»What are you doing here?« she said, and she took Mary by the arm and pulled her away. »What did I tell you?«

»I turned round the wrong corner,« explained Mary. »I didn't know which way to go and I heard some one crying.«

She quite hated Mrs~Medlock at the moment, but she hated her more the next.

»You didn't hear anything of the sort,« said the housekeeper. »You come along back to your own nursery or I'll box your ears.«

And she took her by the arm and half pushed, half pulled her up one passage and down another until she pushed her in at the door of her own room.

»Now,« she said, »you stay where you're told to stay or you'll find yourself locked up. The master had better get you a governess, same as he said he would. You're one that needs some one to look sharp after you. I've got enough to do.«

She went out of the room and slammed the door after her, and Mary went and sat on the hearth-rug, pale with rage. She did not cry, but ground her teeth.

»There was some one crying—there was—there was!« she said to herself.

She had heard it twice now, and sometime she would find out. She had found out a great deal this morning. She felt as if she had been on a long journey, and at any rate she had had something to amuse her all the time, and she had played with the ivory elephants and had seen the gray mouse and its babies in their nest in the velvet cushion.